% !TeX program = xelatex
% !TeX encoding = UTF-8
% 《几何朗兰兹笔记：扩展 Whittaker 范畴》中文译注本
% 底本：D. Gaitsgory, "Notes on Geometric Langlands: The Extended Whittaker Category",
%        preprint, December 2, 2010, 17 pp. (pdfTeX)
% 编译：xelatex -interaction=nonstopmode main.tex（三遍）
\documentclass[UTF8, a4paper, 11pt, openany]{ctexbook}

% ---------- 宏包 ----------
\usepackage{amsmath, amssymb, amsthm, mathtools}
\usepackage{fontspec}
\usepackage{geometry}
\usepackage{fancyhdr}
\usepackage{tikz}
\usepackage{tikz-cd}
\usepackage{enumitem}
\usepackage{booktabs}
\usepackage{xcolor}
\usepackage{hyperref}

% ---------- 字体 ----------
% 本机只有 Noto Serif SC / Noto Sans SC 的可变字体，无独立 Bold/Italic 字面，
% 故用 AutoFakeBold / AutoFakeSlant 合成。
\setCJKmainfont{Noto Serif SC}[AutoFakeBold=1.5, AutoFakeSlant=0.2]
\setCJKsansfont{Noto Sans SC}[AutoFakeBold=1.5]
\setCJKmonofont{Noto Sans SC}
\setmonofont{Consolas}[Scale=0.9]

% ---------- 版式 ----------
\geometry{left=2.8cm, right=2.6cm, top=2.8cm, bottom=2.6cm,
  headsep=0.6cm, footskip=0.9cm}
\setlength{\headheight}{14pt}
\pagestyle{fancy}
\fancyhf{}
\renewcommand{\headrulewidth}{0.4pt}
\fancyhead[LE]{\nouppercase{\leftmark}}
\fancyhead[RO]{\nouppercase{\rightmark}}
\fancyfoot[C]{\thepage}
\fancypagestyle{plain}{%
  \fancyhf{}\renewcommand{\headrulewidth}{0pt}\fancyfoot[C]{\thepage}}

\tolerance=1500
\emergencystretch=3em
\hfuzz=2pt
\raggedbottom

% 节标题左对齐加粗（与原书一致）
\ctexset{section/format=\Large\bfseries\raggedright}

% ---------- 原书编号 ----------
% 原书的编号段与命题/引理共用一套「节.序号」编号（1.1、1.3.1、2.2.2……），
% 编号显式印出、不随排版浮动，故用显式编号宏，交叉引用仍走 \ref。
\makeatletter
\newcounter{oblk}
\newcommand{\ob}[2]{%
  \par\smallskip\refstepcounter{oblk}%
  \def\@currentlabel{#1}\label{#2}%
  \noindent\textbf{#1.}\ \ignorespaces}
\newenvironment{othm}[3]{%
  \par\smallskip\refstepcounter{oblk}%
  \def\@currentlabel{#2}\label{#3}%
  \noindent\textbf{#1 #2.}\ \ignorespaces}{\par\smallskip}
\makeatother

\renewcommand{\proofname}{证明}
\renewcommand{\qedsymbol}{$\square$}

% ---------- 译注 ----------
% 三类：易错预警、含糊点拨、背景直觉。每条附来源等级。
\newenvironment{annote}{%
  \par\vspace{0.3em}\begingroup\small
  \noindent\textcolor{black!50}{\rule[-0.25em]{2.2pt}{1.1em}}\hspace{0.5em}%
  \textsf{\textbf{译注}}\hspace{0.4em}%
}{\par\endgroup\vspace{0.3em}}

% ---------- 常用记号 ----------
\newcommand{\A}{\mathbb{A}}
\newcommand{\Gm}{\mathbb{G}_m}
\newcommand{\Ga}{\mathbb{G}_a}
\newcommand{\Bun}{\operatorname{Bun}}
\newcommand{\LocSys}{\operatorname{LocSys}}
\newcommand{\Gr}{\operatorname{Gr}}
\newcommand{\Hom}{\operatorname{Hom}}
\newcommand{\Spec}{\operatorname{Spec}}
\newcommand{\Rep}{\operatorname{Rep}}
\newcommand{\Loc}{\operatorname{Loc}}
\newcommand{\Eis}{\mathrm{Eis}}
\newcommand{\CT}{\mathrm{CT}}
\newcommand{\Whit}{\mathrm{Whit}}
\newcommand{\Av}{\mathrm{Av}}
\newcommand{\Char}{\mathrm{Char}}
\newcommand{\oblv}{\mathbf{oblv}}
\newcommand{\ind}{\mathbf{ind}}
\newcommand{\coeff}{\mathsf{coeff}}
\newcommand{\Poinc}{\mathsf{Poinc}}
\newcommand{\QCoh}{\mathsf{QCoh}}
\newcommand{\IndCoh}{\mathsf{IndCoh}}
\newcommand{\fP}{\mathfrak{p}}
\newcommand{\fQ}{\mathfrak{q}}
\newcommand{\fI}{\mathfrak{i}}
\newcommand{\fJ}{\mathfrak{j}}
\newcommand{\fR}{\mathfrak{r}}
\newcommand{\bfP}{\bar{\mathfrak{p}}}
\newcommand{\bfQ}{\bar{\mathfrak{q}}}
\newcommand{\bBun}{\overline{\mathrm{Bun}}}
\newcommand{\Ch}{\check}
\newcommand{\dg}{^{\dagger}}
\newcommand{\rt}{^{rat}}
\newcommand{\rd}{^{red}}
\newcommand{\ext}{^{red,ext}}
\newcommand{\st}{^{spec}}
\newcommand{\sN}{^{N}}
\newcommand{\one}{\mathbf{1}}

\theoremstyle{plain}
\newtheorem*{plainthm}{定理}

% ---------- 超链接 ----------
\hypersetup{
  colorlinks=true, linkcolor=blue!70!black, citecolor=green!50!black,
  urlcolor=blue!70!black, bookmarks=true, bookmarksnumbered=true,
  pdftitle={几何朗兰兹笔记：扩展 Whittaker 范畴（中文译注）},
  pdfauthor={D. Gaitsgory}
}

\setcounter{tocdepth}{1}

% 正文不设 \chapter，节号须与原著 1--8 一致，故去掉章节前缀
\renewcommand{\thesection}{\arabic{section}}
\renewcommand{\theHsection}{\arabic{section}}
% 原书的公式编号也是「节.序号」（(4.1)、(7.1)、(8.1)），故公式按节重新编号
\makeatletter
\@addtoreset{equation}{section}
\makeatother
\renewcommand{\theequation}{\thesection.\arabic{equation}}

\begin{document}

\frontmatter
\thispagestyle{empty}
\begin{center}
\vspace*{2.2cm}

{\Large\bfseries 几何朗兰兹笔记}\\[0.35em]
{\Large\bfseries 扩展 Whittaker 范畴}

\vspace{2.2cm}

{\large Dennis Gaitsgory 著}

\vspace{0.6cm}

{\large 中文译注本}

\vspace{2.6cm}

{\normalsize 底本：\emph{Notes on Geometric Langlands:}\\[0.2em]
\emph{The Extended Whittaker Category}\\[0.5em]
二〇一〇年十二月二日，十七页}

\vfill

{\small 工程编号 SY02-44\quad 二〇二六年十一月}
\end{center}
\clearpage

\chapter*{译者说明}
\markboth{译者说明}{}
\addcontentsline{toc}{chapter}{译者说明}

这份译注本译自 D. Gaitsgory 的预印本 \emph{Notes on Geometric Langlands: The Extended Whittaker Category}（2010 年 12 月 2 日，17 页）。它是作者「几何朗兰兹笔记」系列中的一册，该系列以 [GL:…] 为代号互相引用，本文多处结论的证明留在尚未写出的后续册中，作者用「Quasi-」（拟）标出。译文保留这一区分：凡原文标 Quasi- 者，译作「拟定理」「拟命题」等，不因中文读来生硬而抹平。

\medskip
\noindent\textbf{几点体例}

\begin{enumerate}[leftmargin=2.2em, itemsep=0.35em]
\item 原书的节号（1--8）、编号段（1.1、2.1.1……）与定理类编号（1.3.1、2.2.1……）共用一套数字，译文照旧，不另编号。
\item 原书的映射用 Fraktur 体 $\fP,\fQ,\fI,\fJ,\fR$，加横线者为紧化版本（$\bfP,\bfQ$），再加上标 $rat,red,ext,spec$ 层层区分；这些记号在中文排版中极易看错，译文一律原样保留，并在首次出现处加括注。
\item 原文的笔误（如 \texttt{consant}、$\mathrm{Av}_{\mathcal H_{G,H}}$、\texttt{neither zeros not poles}、定理编号误写、§3.2 两处恒等式的函子方向不合）均按「保留原文、另加译注」处理，不径改原文。各条的理由见工程文件 \texttt{work/疑点清单.tsv}。
\item 每条解释性译注都标注来源：\textbf{A} 指原书、作者他文或经典教材，\textbf{B} 指 nLab、Stacks Project、ProofWiki 一类公开辞书（须与 A 交叉），\textbf{I} 指本仓库已交付工程的术语表；凡无外部依据、仅由上下文推断者，标「译者按（上下文推断，无外部依据）」。
\item 术语依工程文件 \texttt{work/glossary.tsv}（冻结版）统一；「既约（reduced）」「系数函子」「Poincaré 级数函子」等取舍的理由见该表与 \texttt{work/通读笔记.md}。
\end{enumerate}

\medskip
\noindent\textbf{原书页与本书节的对照}

\begin{center}
\begin{tabular}{lll}
\toprule
原书页 & 内容 & 本书 \\
\midrule
p.1 & 引言 & 引言 \\
p.1--2 & §1 到 Borel 的有理约化空间 & §1 \\
p.2--4 & §2 施加 $N(\A)$-等变 & §2 \\
p.4--5 & §3 Eisenstein 级数函子 & §3 \\
p.6--8 & §4 谱图景 & §4 \\
p.8--10 & §5 Whittaker 范畴 & §5 \\
p.10--11 & §6 既约 Whittaker 范畴 & §6 \\
p.12--16 & §7 扩展 Whittaker 范畴 & §7 \\
p.16--17 & §8 粘合与谱图景 & §8 \\
p.17 & 参考文献 & 文献 \\
\bottomrule
\end{tabular}
\end{center}

\medskip
\noindent 底本为作者未正式发表的预印本，权利归原作者；本译本仅作内部研读之用，不对外发布。

\tableofcontents

\mainmatter
\chapter*{引言}
\markboth{引言}{}
\addcontentsline{toc}{chapter}{引言}

这一系列笔记的目的是解释与几何朗兰兹对应有关的当前想法。该对应应当把范畴 $D(\Bun_G)$ 与叠 $\LocSys_{\Ch G}$ 上的拟凝聚层范畴联系起来。

本册研究这些猜想性构造所涉及的几何，更确切地说，所涉及的范畴。它们大多是经典自守形式论中某些函数空间的几何类比，例如主系列与 Whittaker 模型。

关于术语的一点说明：我们说「范畴」，指的是 [GL:DG] 所采用约定下的 DG 范畴；极限与余极限都在该文记号下的 $(\infty,1)$-范畴 $\mathrm{DGCat}_{cont}$ 中取。

关于 $\Bun_G$ 或 $\Bun_B$ 这类叠上 D-模范畴的定义，见 [DrGa1]。

\begin{annote}
本册在系列中的位置：§4 的谱侧讨论遵循 [GL:Eis]；§7--§8 的结论供 [GL:funct] 使用，后者据以对 $G=GL_2$ 构造几何朗兰兹函子 $\Psi_G$；§1 的拟定理 1.3.1 与 §5 的拟定理 5.3.3 分别在 [GL:contr]、[Ga2] 中证明。读者若只关心主角「扩展 Whittaker 范畴」，可先读 §6--§7，再回头补 §1--§4 的记号。\par\vspace{0.15em}{\raggedleft 【来源：原书各处引用（A）】\par}
\end{annote}

\section{到 Borel 的有理约化空间}
\markboth{到 Borel 的有理约化空间}{}

众所周知，叠 $\Bun_B$ 不连通。因此，若想用伴随函子对
\[
  \fP_! : D(\Bun_B)\rightleftarrows D(\Bun_G) : \fP^!
\]
来逼近范畴 $D(\Bun_G)$，就会丢失大量信息。

为此我们转向 Drinfeld 紧化 $\bBun_B$ 及其映射
\[
  \bfP : \bBun_B \to \Bun_G .
\]
然而 $\bBun_B$ 另有一个毛病：冗余太多。也就是说，它不只是把 $\Bun_B$ 的各个连通分支彼此粘起来，而是沿着曲线的部分对称化幂把各分支的乘积粘起来。

本节的目标是定义一个「空间」$\Bun\rt_B$，它把 $\Bun_B$ 的所有连通分支真正无冗余地粘在一起。唯一的问题是：这个空间不存在。不过范畴 $D(\Bun\rt_B)$ 存在，而这正是我们要做的事。想法很简单：考察 $D(\bBun_B)$，再用关于适当群胚的等变性把冗余消去。

\ob{1.1}{blk:1.1}\textbf{「空间」。}考虑叠 $\bBun_B$。令 $\mathcal H_{G,B}$ 为群胚
\[
  \mathcal H_{G,B}\subset \bBun_B\underset{\Bun_G}{\times}\bBun_B ,
\]
它由给定 $G$-丛上、在曲线的一般点处相等的广义 $B$-约化对构成。

我们把范畴 $D(\Bun\rt_B)$ 定义为 $D(\bBun_B)$ 关于 $\mathcal H_{G,B}$ 的等变范畴。

\emph{注。}可以证明 $\Bun\rt_B$ 不作为 Artin 叠存在，即不存在一个叠，带有映射 $\Bun_B\to\Bun\rt_B$ 且该映射诱导出遗忘函子 $D(\Bun\rt_B)\to D(\Bun_B)$。\footnote{不存在的证明来自：$D(\Bun\rt_B)$ 上不存在具有预期性质的 $t$-结构。}然而为记号方便，我们仍把它当作这样的叠来处理。

\begin{annote}
「广义 $B$-约化」是 Drinfeld 紧化 $\bBun_B$ 的原始定义中的对象：在每个点处允许 $B$-约化带有极点，极点由一个 $G$-支配余特征标刻画。两点在一般点处相等，等价于说两者的差别由曲线的除子给出；这些差别正是下面 $X^\lambda$ 所参数化的冗余。\par\vspace{0.15em}{\raggedleft 【来源：原书 p.2 §1.1（A）；[GL:contr] 的引言（A）】\par}
\end{annote}

\ob{1.2}{blk:1.2}\textbf{单子描述。}令 $\oblv_B$ 记遗忘函子
\[
  D(\Bun\rt_B)\to D(\bBun_B).
\]
投影
\[
  \bBun_B\xleftarrow{h^\leftarrow}\mathcal H_{G,B}\xrightarrow{h^\rightarrow}\bBun_B
\]
都是 ind-紧合的。因此函子 $\oblv_B$ 有左伴随，记作 $\ind_B$。复合
\[
  \oblv_B\circ\ind_B : D(\bBun_B)\to D(\bBun_B)
\]
等同于 $\Av_{\mathcal H_{G,B}} := h^\rightarrow_*\circ h^!_\leftarrow$。由 Barr-Beck，我们可以把 $D(\Bun\rt_B)$ 看作 $D(\bBun_B)$ 上关于单子 $\Av_{\mathcal H_{G,B}}$ 的模。

\begin{annote}
为何 ind-紧合给出左伴随：对 ind-紧合映射，$!$-拉回与 $*$-推前构成伴随对；等变范畴的遗忘函子是取关于该群胚的「平均」，而平均由 $h^\rightarrow_*\circ h^!_\leftarrow$ 实现，故遗忘函子的左伴随是沿 $h^\rightarrow$ 的 $!$-推前。Barr-Beck 定理则把等变范畴等同于该单子的模范畴。\par\vspace{0.15em}{\raggedleft 【来源：原书 p.2 §1.2（A）；[GL:DG] 关于 D-模六函子形式的约定（A）】\par}
\end{annote}

\ob{1.3}{blk:1.3}\textbf{与 $\Bun_G$ 的关系。}有伴随函子对
\[
  \bfP_! : D(\bBun_B)\rightleftarrows D(\Bun_G) : \bfP^! .
\]
函子 $\bfP^!$ 典范地分解为 $\oblv_B\circ(\fP\rt)^!$，其中 $(\fP\rt)^!$ 是典范定义的函子 $D(\Bun_G)\to D(\Bun\rt_B)$。于是 $(\fP\rt)^!$ 有左伴随，记作 $(\fP\rt)_!$，使得
\[
  \bfP \simeq (\fP\rt)_!\circ\ind_B .
\]
下述将在 [GL:contr] 中证明：

\begin{othm}{拟定理}{1.3.1}{thm:1.3.1}
函子 $(\fP\rt)^!$ 全忠实。
\end{othm}

\begin{annote}
分解 $\bfP^!\simeq\oblv_B\circ(\fP\rt)^!$ 的含义：先在 $\Bun_G$ 上取 $!$-拉回，再沿 $\Bun\rt_B\to\Bun_G$ 的「部分定义」拉回；后者之所以存在，是因为 $\bfP^!$ 的像落在 $\mathcal H_{G,B}$-等变对象里。由此 $(\fP\rt)^!$ 是 $\bfP^!$ 在等变范畴中的「提升」。拟定理 1.3.1 说这一提升不丢失信息——它是本文反复调用的基本事实（§7.5 的 Quasi-Theorem 7.5.3 是其抛物推广）。\par\vspace{0.15em}{\raggedleft 【来源：原书 p.2 §1.3（A）；译者按（上下文推断，无外部依据）】\par}
\end{annote}

\section{施加 \texorpdfstring{$N(\A)$}{N(A)}-等变}
\markboth{施加 \texorpdfstring{$N(\A)$}{N(A)}-等变}{}

范畴 $D(\Bun\rt_B)$ 应当是双商
\[
  B(K)\backslash G(\A)/G(\mathbb O)
\]
上函数空间的几何类比。

现在引入一个范畴，记作 $D(\overline{\Bun}_T)$，它是双商
\[
  T(K)N(\A)\backslash G(\A)/G(\mathbb O)
\]
上函数空间的几何类比。

注意，正如在集合论意义上有
\[
  B(K)\backslash G(\A)/G(\mathbb O)\simeq B(K)\backslash B(\A)/B(\mathbb O),
\]
我们也有
\[
  T(K)N(\A)\backslash G(\A)/G(\mathbb O)\simeq T(K)\backslash T(\A)/T(\mathbb O).
\]
然而，正如范畴 $D(\Bun\rt_B)$ 不等价于 $D(\Bun_B)$，$D(\Bun\rt_T)$ 也不会等同于 $D(\Bun_T)$。

\ob{2.1}{blk:2.1}\textbf{范畴。}叠 $D(\bBun_B)$ 带一个典范群胚的作用，记作 $\mathcal H_N$，它实现 $N(\A)$ 的 adelic 作用。读者可参阅 [Ga1] 或 [Ga2]，那里对此有详细解释。\footnote{原文把「范畴」写成了「叠」：带作用的是范畴 $D(\bBun_B)$，$\mathcal H_N$ 是作用在其上的群胚。译文照录原文，读者按范畴理解即可。}

令 $D(\bBun_T)$ 为 $D(\bBun_B)$ 中 $\mathcal H_N$-等变对象的范畴。由于 $\mathcal H_N$ 投射幂幺，遗忘函子
\[
  D(\bBun_T)\to D(\bBun_B)
\]
全忠实。我们把这个函子记作 $\bfQ^*$。它有右伴随，记作 $\bfQ_*$。

\emph{注。}尽管除 $GL_2$ 的情形外我们不能把 $\bBun_T$ 定义为叠，为记号方便我们仍如此处理。

\ob{2.1.1}{blk:2.1.1}容易看出群胚 $\mathcal H_{G,B}$ 与 $\mathcal H_N$ 的作用交换。因此，考虑 $D(\Bun\rt_B)$ 中 $\mathcal H_N$-等变对象的范畴是合理的；它是 $D(\Bun\rt_B)$ 的全子范畴，我们记作 $D(\Bun\rt_T)$。

等价地，注意作用在 $D(\bBun_B)$ 上的单子 $\Av_{\mathcal H_{G,B}}$ 把 $D(\bBun_T)$ 送到自身。我们可以把 $D(\Bun\rt_T)$ 描述为 $D(\bBun_T)$ 中 $\Av_{\mathcal H_{G,B}}$-模的范畴。

把相应的全忠实嵌入记作
\[
  (\fQ\rt)^* : D(\Bun\rt_T)\hookrightarrow D(\Bun\rt_B),
\]
其右伴随记作 $(\fQ\rt)_*$。

记
\[
  \ind_T : D(\bBun_T)\rightleftarrows D(\Bun\rt_T) : \oblv_T
\]
为所得的伴随函子对。

\ob{2.2}{blk:2.2}\textbf{与 $D(\Bun_T)$ 的关系。}开嵌入
\[
  \fJ_B : \Bun_B\hookrightarrow\bBun_B
\]
诱导开嵌入
\[
  \fJ_T : \Bun_T\hookrightarrow\bBun_T,
\]
以及相应的函子
\[
  (\fJ_T)^! : D(\bBun_T)\to D(\Bun_T).
\]

\begin{othm}{引理}{2.2.1}{thm:2.2.1}
函子 $(\fJ_T)^!$ 有左伴随。
\end{othm}

换言之，引理说的是：$(\fJ_B)^!$ 的（部分定义的）左伴随在全子范畴 $D(\Bun_T)\subset D(\Bun_B)$ 上有定义。理由如下：对象 $(\fJ_B)_!(k_{\Bun_B})\in D(\bBun_B)$（因完整（holonomic）而有定义）关于投影 $\bfQ:\bBun_B\to\bBun_T$ 是 ULA 的。\footnote{这里 $k_Y$ 记概形或叠 $Y$ 上的「常值层」D-模。}

把 $(\fJ_T)^!$ 的左伴随记作 $(\fJ_T)_!$。

\ob{2.2.2}{blk:2.2.2}考虑复合函子
\[
  (\fJ\rt_T)_! := \ind_T\circ(\fJ_T)_! : D(\Bun_T)\to D(\Bun\rt_T),
\]
及其右伴随
\[
  (\fJ\rt_T)^! := \fJ_T^!\circ \oblv_T .
\]
容易看出函子 $(\fJ\rt_T)^!$ 保守。因此范畴 $D(\Bun\rt_T)$ 可以描述为 $D(\Bun_T)$ 上关于单子
\[
  \widetilde\Omega := (\fJ\rt_T)^!\circ(\fJ\rt_T)_!
  \simeq (\fJ_T)^!\circ\oblv_T\circ\ind_T\circ(\fJ_T)_!
\]
的模。（记号 $\widetilde\Omega$ 源自 [GL:Eis] 的记号约定。）

\begin{annote}
$(\fJ\rt_T)^!$ 保守（即它只把同构映到同构）是「单子模范畴」描述的关键：正因为保守，$D(\Bun\rt_T)$ 才能由它在 $D(\Bun_T)$ 上的像与单子 $\widetilde\Omega$ 完全刻画。证明的思路是：$(\fJ\rt_T)^!$ 是 $\bfQ^*$ 与 $\fJ_T^!$ 的复合，前者在等变范畴上是「取底层对象」，后者在 $\Bun_T$ 上是限制，二者在相应范畴上都保守。\par\vspace{0.15em}{\raggedleft 【来源：原书 p.4 §2.2.2（A）；译者按（上下文推断，无外部依据）】\par}
\end{annote}

\ob{2.3}{blk:2.3}\textbf{单子的描述。}单子 $\widetilde\Omega$ 可以显式描述。对元素 $\lambda\in\Lambda$，令
\[
  \fJ_B^\lambda : X^\lambda\times\Bun_B\to\bBun_B
\]
为相应分层的局部闭嵌入。我们用 $\fJ_T^\lambda$ 记号地表示「局部闭嵌入」$X^\lambda\times\Bun_T\to\bBun_T$；我们的意思只是：存在函子
\[
  (\fJ_T^\lambda)^! : D(\bBun_T)\to D(X^\lambda\times\Bun_T).
\]
如同引理 2.2.1，函子 $(\fJ_T^\lambda)^!$ 有左伴随，我们记作 $(\fJ_T^\lambda)_!$。

考虑函子
\[
  \widetilde\Omega^\lambda : D(\Bun_T)\to D(\Bun_T),
\]
它由
\[
  \bigl(H(X^\lambda,-)\boxtimes \mathrm{Id}_{D(\Bun_T)}\bigr)\circ(\fJ_T^\lambda)^!\circ(\fJ_T)_!
\]
给出。

由构造，函子
\[
  \bigoplus_{\lambda\in\Lambda^{pos}}\widetilde\Omega^\lambda
\]
带有自然的单子结构。

\begin{othm}{引理}{2.3.1}{thm:2.3.1}
上述单子 $\bigoplus_{\lambda\in\Lambda^{pos}}\widetilde\Omega^\lambda$ 典范同构于 $\widetilde\Omega$。
\end{othm}

引理来自同构
\[
  \Bun_B\underset{\bBun_B}{\times}\mathcal H_{G,B}\simeq\bigsqcup_\lambda X^\lambda\times\Bun_B,
\]
其中「第二投影」$\Bun_B\times_{\bBun_B}\mathcal H_{G,B}\to\bBun_B$ 由映射 $\bigsqcup_\lambda\fJ_B^\lambda$ 给出。

\begin{annote}
为何只对 $\lambda\in\Lambda^{pos}$ 求和：$\mathcal H_{G,B}$ 的两点只在一般点处相等，从第一个约化变到第二个约化所需的除子必为非负组合，故参数 $\lambda$ 落在正半群 $\Lambda^{pos}$ 中。$X^\lambda$ 是曲线 $X$ 的相应部分对称化幂；当 $\lambda=0$ 时 $X^0=\mathrm{pt}$，$\widetilde\Omega^0$ 就是恒等函子。\par\vspace{0.15em}{\raggedleft 【来源：原书 §2.3 与 [GL:contr] 中 $\Bun_B\times_{\bBun_B}\mathcal H_{G,B}$ 的分解（A）】\par}
\end{annote}

\section{Eisenstein 级数函子}
\markboth{Eisenstein 级数函子}{}

在经典情形，Eisenstein 级数与常数项是在
\[
  G(K)\backslash G(\A)/G(\mathbb O)
  \quad\text{与}\quad
  T(K)N(\A)\backslash G(\A)/G(\mathbb O)
\]
上的函数空间之间起作用的算子，其定义是沿图
\[
  G(K)\backslash G(\A)/G(\mathbb O)\leftarrow B(K)\backslash G(\A)/G(\mathbb O)
  \to T(K)N(\A)\backslash G(\A)/G(\mathbb O)
\]
作拉回—推前。下面我们定义这一手续的几何类比。

\ob{3.1}{blk:3.1}\textbf{紧化与非紧化的几何 Eisenstein 级数。}回忆紧化 Eisenstein 级数函子的定义：
\[
  \overline{\Eis} := \bfP_!\circ\bfQ^* : D(\bBun_T)\to D(\Bun_G).
\]
它的左伴随是「紧化常数项」函子：
\[
  \overline{\CT} := \bfQ_*\circ\bfP^! : D(\Bun_G)\to D(\bBun_T).
\]

复合函子 $\Eis_! := \overline{\Eis}\circ(\fJ_T)_!$ 是通常的（非紧化）Eisenstein 级数函子。它的左伴随 $\CT$ 同构于 $\fJ_T^{\bar*}\circ\overline{\CT}$。

\begin{annote}
$\fJ_T^{\bar*}$ 是原文未加定义的记号。由类型与伴随关系推断，它应读作 $\fJ_T$ 的「重新正规化的 $\ast$-拉回」，即 $(\fJ_T)_!$ 的部分定义左伴随：$\Eis_! = \overline{\Eis}\circ(\fJ_T)_!$ 的左伴随，正是 $(\fJ_T)_!$ 的左伴随与 $\overline{\CT}$ 的复合。译文照录原记号，不代读者改写。\par\vspace{0.15em}{\raggedleft 【来源：原书 p.5 §3.1（A）；译者按（上下文推断，无外部依据）】\par}
\end{annote}

\ob{3.2}{blk:3.2}\textbf{有理 Eisenstein 级数。}记 $\Eis\rt$ 为函子
\[
  \fP\rt_!\circ(\fQ\rt)^* : D(\Bun\rt_T)\to D(\Bun_G),
\]
并记 $\CT\rt$ 为其右伴随。

我们把这两个函子看作本节开头提到的 Eisenstein 算子与常数项算子的几何对应物。

同义反复地，我们有：
\[
  \overline{\Eis}\simeq\Eis\rt\circ\ind_T
  \quad\text{与}\quad
  \CT\rt\simeq\oblv_T\circ\overline{\CT},
\]
\[
  \Eis_!\simeq\Eis\rt\circ(\fJ\rt_T)_!
  \quad\text{与}\quad
  \CT\simeq(\fJ\rt_T)^!\circ\mathcal T_r .
\]

\begin{annote}
上面两行恒等式中的第二式与函子方向不合：$\overline{\CT}$ 取值于 $D(\bBun_T)$，而 $\oblv_T$ 的定义域是 $D(\Bun\rt_T)$；$(\fJ\rt_T)^!$ 亦然，它只能右接取值于 $D(\Bun\rt_T)$ 的函子。按伴随关系，$\overline{\Eis}\simeq\Eis\rt\circ\ind_T$ 取右伴随给出 $\overline{\CT}\simeq\oblv_T\circ\CT\rt$，而 $\Eis_!\simeq\Eis\rt\circ(\fJ\rt_T)_!$ 取左伴随给出 $\CT\simeq(\fJ\rt_T)^!\circ\CT\rt$。故原文应是把 $\CT\rt$ 与 $\overline{\CT}$ 写反，并把 $\CT\rt$ 误排成了未定义的 $\mathcal T_r$。译文照录原文，以便与底本对读；读者按上述两式理解即可。\par\vspace{0.15em}{\raggedleft 【来源：原书 p.5 §3.2（A）；译者按（类型检查与伴随性验算）】\par}
\end{annote}

\ob{3.3}{blk:3.3}\textbf{常数项与 Eisenstein 级数。}在分析 $D(\Bun_G)$ 时，一个基本的角色由「计算 Eisenstein 级数之间的 Hom」这一能力扮演。由伴随性，我们有
\[
  \Hom_{D(\Bun_G)}\bigl(\Eis_!(\mathcal M_1),\Eis_!(\mathcal M_2)\bigr)
  \simeq \Hom_{D(\Bun_T)}\bigl(\mathcal M_1,\CT\circ\Eis_!(\mathcal M_2)\bigr).
\]
也就是说，我们需要理解单子
\[
  \CT\circ\Eis_! : D(\Bun_T)\to D(\Bun_T)
\]
的结构。

由于函子 $\Eis_!$ 可以分解为 $\Eis\rt\circ(\fJ\rt_T)_!$，我们有单子的典范映射
\[
  \widetilde\Omega\to\CT\circ\Eis_! .
\]

下述命题来自用 Zastava 空间对单子 $\widetilde\Omega$ 的计算，以及引理 2.3.1：

\begin{othm}{命题}{3.3.1}{prop:3.3.1}
函子 $\CT\circ\Eis_!$ 带有按 Weyl 群元素 $w\in W$ 编号的典范滤过，使得：
\begin{enumerate}[label=(\alph*), leftmargin=2.4em, itemsep=0.25em]
\item 首项 $(\CT\circ\Eis_!)_{unit}$ 带有单子结构，且作为单子典范同构于 $\widetilde\Omega$；
\item 末项 $(\CT\circ\Eis_!)_{w_0}$ 典范同构于 $w_0$ 在 $D(\Bun_T)$ 上的作用函子。
\end{enumerate}
\end{othm}

\emph{注。}命题的 (a) 可以这样理解：单子 $\widetilde\Omega$ 是单子 $\CT\circ\Eis_!$ 在关于 $\Bun_T$ 的谱意义上「局部」的最佳逼近。

还要指出：在经典的自守函数论中，单子 $\widetilde\Omega$ 的类比物是 $T$ 所对应的 adelic 商上函数空间的一个自同态，它在某个 Grossen 特征（Größencharakter）上的取值，等于（所有正根之积对应的）阿贝尔 $L$ 函数与其移位变元之比的商。

\begin{annote}
两处需要点破。其一，「按 $w\in W$ 的滤过」是指把 $\CT\circ\Eis_!$ 写成归次为 $w$ 的逐次商之链，首项（归次为 $unit$）是单子 $\widetilde\Omega$，末项（归次为 $w_0$）是 $w_0$ 的作用；这正是 §4.5 在谱侧对 $(\fP\st_N)^!\circ\fP\st_*$ 所证的同型结论。其二，经典类比里的「比值」应读作 $L(\chi,s)/L(\chi,s+?)$ 形状的商：正根之积给出 $\prod_{\alpha>0}(1-\chi(\alpha)t)$ 型的因子，移位变元来自 Weyl 向量 $\rho$。\par\vspace{0.15em}{\raggedleft 【来源：原书 p.5 §3.3（A）；经典自守形式论中 Eisenstein 级数的常数项公式（A）】\par}
\end{annote}

\section{谱图景}
\markboth{谱图景}{}

\ob{4.1}{blk:4.1}\textbf{谱 Eisenstein 级数。}考虑叠的映射
\[
  \LocSys_{\Ch G}\xleftarrow{\fP\st}\LocSys_{\Ch B}\xrightarrow{\fQ\st}\LocSys_{\Ch T}.
\]
谱 Eisenstein 级数函子
\[
  \Eis\st : \QCoh(\LocSys_{\Ch T})\to \QCoh_N(\LocSys_{\Ch G})
\]
定义为复合 $(\fP\st)_*\circ(\fQ\st)^*$。这里 $\QCoh_N(\LocSys_{\Ch G})$ 是 $\QCoh(\LocSys_{\Ch G})$ 的放大，由 [GL:formulation] 引入。

把 $\Eis\st$ 的右伴随记作 $\CT\st$，称之为谱常数项函子。\footnote{为使 $\CT\st$ 作为保余极限的函子存在，需要 $\Eis\st$ 把紧对象送到紧对象。这正是引入修改 $\QCoh(\LocSys_{\Ch G})\to\QCoh_N(\LocSys_{\Ch G})$ 的原因：我们需要把 $\QCoh(\LocSys_{\Ch G})$ 的紧对象从「只是完美复形」扩大到包含 $\Eis\st$ 的像中的那些凝聚复形。}

几何朗兰兹等价所应具备的性质之一是：它把函子 $\Eis_!$ 与（按 $\rho(\omega_X)$ 平移后的）$\Eis\st$ 相缠，也就是说，下图必须交换：
\begin{equation}
\begin{tikzcd}[row sep=3.4em, column sep=3.6em, ampersand replacement=\&]
  D(\Bun_T) \& \QCoh(\LocSys_{\Ch T})
    \arrow[l, "-\rho(\omega_X)\circ\Psi_T"'] \\
  D(\Bun_G) \& \QCoh_N(\LocSys_{\Ch G})
    \arrow[l, "\Psi_G"']
    \arrow[from=1-1, to=2-1, "\Eis_!"]
    \arrow[from=1-2, to=2-2, "\Eis\st"']
\end{tikzcd}
\label{eq:4.1}
\end{equation}
这里 $\Psi_G$ 记 $G$ 的猜想性几何朗兰兹函子，而 $\Psi_T$ 是实现 $T$ 的几何朗兰兹对应的 Fourier-Mukai 变换。

本节的目标——本节遵循 [GL:Eis]——是解释范畴 $D(\Bun\rt_T)$ 与函子 $\Eis\rt$ 的谱意义。

\ob{4.2}{blk:4.2}\textbf{带纵向联络的拟凝聚层。}再次考虑映射
\[
  \fP\st : \LocSys_{\Ch B}\to\LocSys_{\Ch G},
\]
并令 $T_{\LocSys_{\Ch B}/\LocSys_{\Ch G}}$ 为 $\LocSys_{\Ch B}$ 上关于 $\fP\st$ 的纵向量场代数胚。

考虑范畴
\[
  T_{\LocSys_{\Ch B}/\LocSys_{\Ch G}}\text{-}\mathrm{mod}^r,
\]
即 $\QCoh(\LocSys_{\Ch B})$ 中的右 $T_{\LocSys_{\Ch B}/\LocSys_{\Ch G}}$-模构成的范畴。

泛包络代数 $U\bigl(T_{\LocSys_{\Ch B}/\LocSys_{\Ch G}}\bigr)$ 可以看作 $\QCoh(\LocSys_{\Ch B})$ 上的一个单子
\[
  \QCoh(\LocSys_{\Ch B})\to\QCoh(\LocSys_{\Ch B}),
\]
它等于遗忘函子
\[
  T_{\LocSys_{\Ch B}/\LocSys_{\Ch G}}\text{-}\mathrm{mod}^r\to\QCoh(\LocSys_{\Ch B})
\]
与其左伴随（即自由函子）的复合。

\begin{annote}
原文此句末作「i.e., the forgetful functor」，与前半句自相矛盾：遗忘函子的左伴随是自由函子 $\mathcal F\mapsto U(T)\otimes\mathcal F$，单子 $U(T)\otimes-$ 是「遗忘 $\circ$ 自由」，而不是遗忘函子本身。译文按「自由函子」译出，原文照录于工程文件疑点清单。\par\vspace{0.15em}{\raggedleft 【来源：原书 p.6 §4.2（A）；译者按（单子与伴随函子的一般理论）】\par}
\end{annote}

\ob{4.3}{blk:4.3}\textbf{精细谱 Eisenstein 级数。}函子
\[
  \fP\st_* : \QCoh(\LocSys_{\Ch B})\to\QCoh_N(\LocSys_{\Ch G})
\]
典范地分解为
\[
  \QCoh(\LocSys_{\Ch B})\xrightarrow{U(T_{\LocSys_{\Ch B}/\LocSys_{\Ch G}})}
  T_{\LocSys_{\Ch B}/\LocSys_{\Ch G}}\text{-}\mathrm{mod}^r
  \xrightarrow{\fP\st_{*,DR}}
  \QCoh_N(\LocSys_{\Ch G}).
\]
由于稍后将要解释的理由，我们把上述函子 $\fP\st_{*,DR}$ 称作「精细谱 Eisenstein 级数」，并记作 $\Eis^{\st,ref}$。

\ob{4.4}{blk:4.4}\textbf{精细谱常数项。}现在考虑 $\fP\st_*$ 的右伴随，我们把它记作 $(\fP\st_N)^!$。注意 $(\fP\st_N)^!$ 是如下复合：
\[
  \QCoh_N(\LocSys_{\Ch G})\leftarrow \IndCoh(\LocSys_{\Ch G})
  \xleftarrow{(\fP\st_N)^!}\IndCoh(\LocSys_{\Ch B})
  \to \QCoh(\LocSys_{\Ch B}),
\]
其中最后一个箭头是包含函子
\[
  \QCoh(\LocSys_{\Ch B})\hookrightarrow\IndCoh(\LocSys_{\Ch B})
\]
的右伴随。

函子 $(\fP\st_N)^!$ 典范地分解为
\[
  \QCoh_N(\LocSys_{\Ch G})\to
  T_{\LocSys_{\Ch B}/\LocSys_{\Ch G}}\text{-}\mathrm{mod}^r\to
  \QCoh(\LocSys_{\Ch B}),
\]
其中第二个箭头是遗忘函子。

我们把所得的函子 $\QCoh_N(\LocSys_{\Ch G})\to T_{\LocSys_{\Ch B}/\LocSys_{\Ch G}}\text{-}\mathrm{mod}^r$ 记作 $\CT^{\st,ref}$，并称之为「精细谱常数项」函子。函子对 $(\Eis^{\st,ref},\CT^{\st,ref})$ 自然地构成伴随对。

\ob{4.5}{blk:4.5}\textbf{谱侧的 Eisenstein 级数之常数项。}现在考虑单子
\[
  (\fP\st_N)^!\circ\fP\st_* : \QCoh(\LocSys_{\Ch B})\to\QCoh(\LocSys_{\Ch B}),
\]
它控制
\[
  \Hom_{\QCoh_N(\LocSys_{\Ch G})}\bigl(\fP\st_*(\mathcal F_1),\fP\st_*(\mathcal F_2)\bigr),
  \qquad \mathcal F_1,\mathcal F_2\in\QCoh(\LocSys_{\Ch B}).
\]

几乎是同义反复地，我们有：

\begin{othm}{引理}{4.5.1}{thm:4.5.1}
函子 $(\fP\st_N)^!\circ\fP\st_*$ 带有按元素 $w\in W$ 编号的典范滤过，使得：
\begin{enumerate}[label=(\alph*), leftmargin=2.4em, itemsep=0.25em]
\item 首项 $\bigl((\fP\st_N)^!\circ\fP\st_*\bigr)_{unit}$ 带有单子结构，且作为单子典范同构于 $U\bigl(T_{\LocSys_{\Ch B}/\LocSys_{\Ch G}}\bigr)$；
\item 末项 $\bigl((\fP\st_N)^!\circ\fP\st_*\bigr)_{w_0}$ 典范同构于复合
\[
  \QCoh(\LocSys_{\Ch B})\xrightarrow{\iota^*}\QCoh(\LocSys_{\Ch T})
  \xrightarrow{w_0}\QCoh(\LocSys_{\Ch T})
  \xrightarrow{\iota_*}\QCoh(\LocSys_{\Ch B}),
\]
其中 $\iota$ 是自然映射 $\QCoh(\LocSys_{\Ch T})\to\QCoh(\LocSys_{\Ch B})$。
\end{enumerate}
\end{othm}

\begin{annote}
引理 4.5.1 与命题 3.3.1 是同一现象的谱侧与几何侧：两侧的滤过都按 $w\in W$ 编号，首项都是「单子」（分别是 $U(T)$ 与 $\widetilde\Omega$），末项都是 $w_0$ 的作用。§4.6 的拟定理正是断言这两套滤过在 $\Psi_T$ 下逐项相合。\par\vspace{0.15em}{\raggedleft 【来源：原书 p.5 命题 3.3.1 与 p.7 引理 4.5.1 的对照（A）】\par}
\end{annote}

\ob{4.6}{blk:4.6}\textbf{几何图景与谱图景的关系。}回忆 $\fQ\st$ 记投影 $\LocSys_{\Ch B}\to\LocSys_{\Ch T}$。令 $\Psi_T$ 记 Fourier-Mukai 变换
\[
  D(\Bun_T)\simeq\QCoh(\LocSys_{\Ch T}).
\]

[GL:Eis] 的主旨如下：

\begin{othm}{拟定理}{4.6.1}{thm:4.6.1}
存在范畴的典范等价
\[
  \Psi_B : T_{\LocSys_{\Ch B}/\LocSys_{\Ch G}}\text{-}\mathrm{mod}^r
  \xrightarrow{\ \sim\ } D(\Bun\rt_T),
\]
它使下图交换：
\[
\begin{tikzcd}[row sep=3.6em, column sep=3.2em, ampersand replacement=\&]
  D(\Bun\rt_T) \&[-2em] T_{\LocSys_{\Ch B}/\LocSys_{\Ch G}}\text{-}\mathrm{mod}^r
    \arrow[l, "\Psi_B"'] \\
  D(\Bun_T) \& \QCoh(\LocSys_{\Ch T})
    \arrow[l, "-\rho(\omega_X)\circ\Psi_T"']
    \arrow[from=1-1, to=2-1, "(\fJ\rt_T)_!"]
    \arrow[from=1-2, to=2-2, "U(T_{\LocSys_{\Ch B}/\LocSys_{\Ch G}})\circ(\fQ\st)^*"]
\end{tikzcd}
\]
其中 $\rho(\omega_X)$ 是按点 $\rho(\omega_X)\in\Bun_T$ 平移的函子。
\end{othm}

注意，定理 4.6.1 可以改述如下。考虑单子
\[
  (\fQ\st)_*\circ U\bigl(T_{\LocSys_{\Ch B}/\LocSys_{\Ch G}}\bigr)\circ(\fQ\st)^*
  : \QCoh(\LocSys_{\Ch T})\to\QCoh(\LocSys_{\Ch T}),
\]
以及
\[
  \rho(\omega_X)\circ(\fJ\rt_T)_!\circ(\fJ\rt_T)^!\circ-\rho(\omega_X)
  =\rho(\omega_X)\circ\widetilde\Omega\circ-\rho(\omega_X)
  : D(\Bun_T)\to D(\Bun_T).
\]
断言是：$\Psi_T$ 把这两个单子相缠。

\begin{annote}
上述两图与两段改述里的 $\rho(\omega_X)$ 都指「按点平移」的函子，而不是乘以线丛。它与 $\widetilde\Omega$ 的复合放在 $\Psi_T$ 两侧，是为了把 §2.2.2 中「在 $D(\Bun_T)$ 上、未平移」的单子搬到与谱侧 $U(T)$ 对齐的位置；这正是几何朗兰兹里「用 $\rho$ 作中心平移」的常见做法。\par\vspace{0.15em}{\raggedleft 【来源：原书 p.8 §4.6（A）；[GL:Eis] 的做法（A）】\par}
\end{annote}

\ob{4.7}{blk:4.7}\textbf{小结。}于是，若我们已有满足 (4.1) 的函子 $\Psi_G$，则下述图也必须交换：
\begin{equation}
\begin{tikzcd}[row sep=3.6em, column sep=3.6em, ampersand replacement=\&]
  D(\Bun\rt_T) \& T_{\LocSys_{\Ch B}/\LocSys_{\Ch G}}\text{-}\mathrm{mod}^r
    \arrow[l, "-\rho(\omega_X)\circ\Psi_T"'] \\
  D(\Bun_G) \& \QCoh_N(\LocSys_{\Ch G})
    \arrow[l, "\Psi_G"']
    \arrow[from=1-1, to=2-1, "\Eis\rt"]
    \arrow[from=1-2, to=2-2, "\Eis^{\st,ref}"']
\end{tikzcd}
\label{eq:4.2}
\end{equation}

\begin{annote}
(4.2) 相对 (4.1) 的进步在于：把 (4.1) 中「谱侧的 $\Eis\st$」换成了「先 $U(T)$、再 $\Eis^{\st,ref}$」的分解，于是左边也相应地把 $D(\Bun_T)$ 换成 $D(\Bun\rt_T)$、把 $\Eis_!$ 换成 $\Eis\rt$。这一步之所以可能，靠的正是拟定理 4.6.1：它把 $T\text{-}\mathrm{mod}^r$ 与 $D(\Bun\rt_T)$ 等同起来。\par\vspace{0.15em}{\raggedleft 【来源：原书 p.8 §4.7（A）】\par}
\end{annote}

\section{Whittaker 范畴}
\markboth{Whittaker 范畴}{}

延续经典的类比，我们现在想构造一个范畴，它是
\[
  N(K)\backslash G(\A)/G(\mathbb O)
\]
上函数空间的几何对应物。这个范畴将记作 $D(\Bun\rt_N)$。此外，我们将定义一个满子范畴
\[
  \Whit(G)_{glob}\subset D(\Bun\rt_N),
\]
它是 $N(K)\backslash G(\A)/G(\mathbb O)$ 上函数空间中「对 $N(\A)$ 的一个固定非退化特征作等变」所得子空间的几何对应物。

这些范畴的构造用到 Ran 空间上的范畴这一形式，见 [GL:Ran]。

\ob{5.1}{blk:5.1}\textbf{空间。}考虑如下 ind-叠，记作 $\Bun_N^{polar}$，它生活在 $\mathrm{Ran}_X$ 之上：对有限集 $I$，令 $(\Bun_N^{polar})_{X^I}$ 分类五元组 $(\underline x,P_G,P_T,\kappa,\gamma)$，其中 $\underline x\in X^I$，$(P_G,P_T,\kappa)$ 是允许在 $\underline x$ 处有极点的 Plücker 数据，$\gamma$ 是 $T$-丛的同构
\[
  P_T\simeq\rho(\omega_X).
\]

对应
\[
  I\mapsto D\bigl((\Bun_N^{polar})_{X^I}\bigr)
\]
自然地是 Ran 空间上的一个范畴，我们记作 $D\bigl((\Bun_N^{polar})_{X^{fSet}}\bigr)$。它自然地带有单位结构与增广。我们把所得的范畴 $D\bigl((\Bun_N^{polar})_{\mathrm{Ran}_X,un}\bigr)$ 记作 $D(\Bun\rt_N)$。

\ob{5.2}{blk:5.2}\textbf{Whittaker 子范畴。}对每个 $I$，存在某个幂幺群胚 $(\mathcal H_N)_{X^I}$ 作用在 $(\Bun_N^{polar})_{X^I}$ 上（见 [Ga2]；事实上 $(\mathcal H_N)_{X^I}$ 是作用于 $\Bun_B$ 的 $\mathcal H_N$ 的一个版本），它带有一个特征
\[
  \chi : (\mathcal H_N)_{X^I}\to\Ga .
\]

我们把 $\Whit(G)_{X^I}$ 定义为相应的等变范畴
\[
  \Whit(G)_{X^I} := D\bigl((\Bun_N^{polar})_{X^I}\bigr)^{((\mathcal H_N)_{X^I},\chi)}.
\]

$(\mathcal H_N)_{X^I}$ 的幂幺性蕴含遗忘函子
\[
  \oblv_{\Whit,X^I} : \Whit(G)_{X^I}\to D\bigl((\Bun_N^{polar})_{X^I}\bigr)
\]
全忠实。该函子有右伴随，记作 $\ind_{\Whit,X^I}$。

\ob{5.2.1}{blk:5.2.1}我们把对应 $I\mapsto\Whit(G)_{X^I}$ 看作 Ran 空间上的一个范畴，简记作 $\Whit(G)$。它自然地是单位的、带增广的。我们将考虑相应的范畴
\[
  \Whit(G)_{\mathrm{Ran}_X}
  \quad\text{与}\quad
  \Whit(G)_{\mathrm{Ran}_X,un}.
\]
我们还将使用记号 $\Whit(G)_{glob}:=\Whit(G)_{\mathrm{Ran}_X,un}$，以及
\[
  \ind_{\Whit,glob} : D(\Bun\rt_N)\rightleftarrows \Whit(G)_{glob} : \oblv_{\Whit,glob}
\]
记相应的伴随函子对。

注意，范畴 $\Whit(G)_\emptyset$（定义为 $D(\Bun_N)$ 中关于相应群胚的等变对象范畴）典范等价于 $\mathrm{Vect}$。我们把 $\Whit(G)_\emptyset\subset\Whit(G)_{glob}$ 的相应单位对象记作 $\one_{\Whit(G)_{glob}}$。

\ob{5.3}{blk:5.3}\textbf{Casselman-Shalika 公式。}考虑由对称幺半范畴 $\Rep(\Ch G)$ 在 Ran 空间上搭起来的范畴；我们把它记作 $\Rep(\Ch G,X^{fSet})$。

\ob{5.3.1}{blk:5.3.1}相应的范畴 $\Rep(\Ch G,X^I)$ 可以如下显式描述。把 $\Ch G$ 看作 $X$ 上常值的群 D-概形，并令
\[
  \Ch G_{X^{fSet}} : I\mapsto \Ch G_{X^I}
\]
为 Ran 空间上相应的分解群 D-概形。

我们有：$\Rep(\Ch G,X^I)$ 是 $\Ch G_{X^I}$ 在 $X^I$ 上的 D-模（与联络相容者）上的表示范畴。

\ob{5.3.2}{blk:5.3.2}下述与 Casselman-Shalika 公式本质等价：

\begin{othm}{拟定理}{5.3.3}{thm:5.3.3}
存在 $\mathrm{Ran}_X$ 上范畴的自然等价
\[
  \Whit(G)\simeq\Rep(\Ch G,X^{fSet}).
\]
\end{othm}

\begin{annote}
「与 Casselman-Shalika 公式本质等价」是说：经典公式把 Whittaker 函数在某个群元素上的取值写成 $\check G$ 的表示的特征标（Schur 函数）在相应参数处的值；几何侧把「Whittaker 函数空间」换成范畴 $\Whit(G)$，把「$\check G$ 的表示」换成 $\Rep(\Ch G,X^{fSet})$，而「取值」这一操作则升级为范畴等价。\par\vspace{0.15em}{\raggedleft 【来源：原书 p.9 §5.3.2（A）；Casselman-Shalika 公式的经典表述（A）】\par}
\end{annote}

\ob{5.3.4}{blk:5.3.4}由上述定理，我们推出等价
\[
  \Whit(G)_{\mathrm{Ran}_X}\simeq\Rep(\Ch G,\mathrm{Ran}_X)
  \quad\text{与}\quad
  \Whit(G)_{\mathrm{Ran}_X,un}\simeq\Rep(\Ch G,\mathrm{Ran}_X,un).
\]
这里我们把 $\Rep(\Ch G,X^{fSet})$ 看作单位的、带增广的，依据是分解（群）D-概形 $I\mapsto\Ch G_{X^I}$ 上的余单位结构。

\ob{5.3.5}{blk:5.3.5}\emph{与 $\QCoh(\LocSys_{\Ch G})$ 的关系。}范畴 $\Rep(\Ch G,\mathrm{Ran}_X)$ 对我们的意义由下述事实说明。回忆有伴随函子对
\[
  \Loc_{\Ch G,spec} : \Rep(\Ch G,\mathrm{Ran}_X)\rightleftarrows \QCoh(\LocSys_{\Ch G}) : \Gamma_{\Ch G,spec}.
\]
我们有：

\begin{othm}{定理}{5.3.6}{thm:5.3.6}
函子 $\Gamma_{\Ch G,spec}$ 全忠实。
\end{othm}

\ob{5.3.7}{blk:5.3.7}容易看出，函子 $\Gamma_{\Ch G,spec}$ 的像落在满子范畴
\[
  \Rep(\Ch G,\mathrm{Ran}_X,un)\subset\Rep(\Ch G,\mathrm{Ran}_X)
\]
中。因此，函子 $\Loc_{\Ch G,spec}$ 经过局部化
\[
  \Rep(\Ch G,\mathrm{Ran}_X)\to\Rep(\Ch G,\mathrm{Ran}_X,un)
\]
分解。

\begin{annote}
定理 5.3.6 与 5.3.7 合起来给出 §7 反复使用的「谱侧落脚点」：$\Whit(G)_{glob}\simeq\Rep(\Ch G,\mathrm{Ran}_X,un)$，而后者经 $\Gamma_{\Ch G,spec}$ 全忠实地坐在 $\QCoh(\LocSys_{\Ch G})$ 里。于是关于 Whittaker 范畴的任何陈述，都可以翻译成谱侧关于拟凝聚层的陈述。\par\vspace{0.15em}{\raggedleft 【来源：原书 p.10 §5.3.5--5.3.7（A）】\par}
\end{annote}

\section{既约 Whittaker 范畴}
\markboth{既约 Whittaker 范畴}{}

从现在起，我们假设 $G$ 的中心（记作 $Z_G$）连通。

为了考虑扩展 Whittaker 范畴，我们需要把上面定义的 $\Whit(G)_{glob}$ 换成它的一个满子范畴，记作 $\Whit(G)\rd_{glob}$。

下面是函数论上的类比。如前所述，$D(\Bun\rt_N)$ 是 $N(K)\backslash G(\A)/G(\mathbb O)$ 上函数空间的类比。我们将引入一个范畴，记作 $D(\Bun^{\mathrm{rat},\mathrm{red}}_N)$，它是
\[
  Z_G(K)N(K)\backslash G(\A)/G(\mathbb O)
\]
上函数空间的几何类比。范畴 $\Whit(G)\rd_{glob}$ 将是 $D(\Bun^{\mathrm{rat},\mathrm{red}}_N)$ 中施加了 Whittaker 型条件所得的满子范畴。

\ob{6.1}{blk:6.1}\textbf{范畴。}令 $\overline{\Bun}\rd_N$ 记分类四元组 $(P_G,P_T,\kappa,\gamma\rd)$ 的叠，其中 $(P_G,P_T,\kappa)$ 是 $\Bun_B$ 的一个点，而 $\gamma\rd$ 是下述资料：对 $G$ 的每个单根 $\Ch\alpha_i$，给定线丛的一个非零映射
\[
  \gamma_i\rd : \Ch\alpha_i(P_T)\to\omega_X .
\]

上述叠带有一个群胚的作用，该群胚记作 $\mathcal H_{G,N}$，它实际上是从 $\Bun_B$ 上的群胚 $\mathcal H_{G,B}$ 拉回而来的：只要存在曲线一般点处定义的同构 $P_T\simeq P_T^1$，使 $\kappa$ 与 $\kappa^1$、$\gamma\rd$ 与 $\gamma^{\mathrm{red},1}$ 相配，它就把两点 $(P_G,P_T,\kappa,\gamma)$ 与 $(P_G,P_T^1,\kappa^1,\gamma^1)$ 等同起来。

我们令 $D(\Bun^{\mathrm{rat},\mathrm{red}}_N)$ 记所得的范畴 $D(\overline{\Bun}\rd_N)^{\mathcal H_{G,N}}$。

与 $D(\Bun\rt_B)$ 的情形一样，遗忘函子
\[
  \oblv\rd_N : D(\Bun^{\mathrm{rat},\mathrm{red}}_N)\to D(\overline{\Bun}\rd_N)
\]
有伴随，记作 $\ind\rd_N$，\footnote{原文作「右伴随」。但紧接着的单子陈述 $\oblv\rd_N\circ\ind\rd_N$ 要求 $\ind\rd_N$ 是左伴随（与 §1.2 的 $\oblv_B\dashv\ind_B$ 一致）；故此处「右」字应为笔误。译文照录原文。}而单子
\[
  \oblv\rd_N\circ\ind\rd_N : D(\overline{\Bun}\rd_N)\to D(\overline{\Bun}\rd_N)
\]
由关于 $\mathcal H_{G,N}$ 取平均给出，即函子 $\Av_{\mathcal H_{G,H}}$。\footnote{原文此处的下标是 $\mathcal H_{G,H}$，全书他处并无 $\mathcal H_{G,H}$；由上下文，应为 $\mathcal H_{G,N}$。译文照录原文。}

\ob{6.2}{blk:6.2}\textbf{Whittaker 条件。}略微滥用记号，我们把从 $\Bun_B$ 拉回到 $\overline{\Bun}\rd_N$ 的同名群胚仍记作 $\mathcal H_N$。注意，$\gamma\rd$ 的资料现在定义出典范特征
\[
  \chi : \mathcal H_N\to\Ga .
\]
由于 $\mathcal H_N$ 与 $\mathcal H_{G,N}$ 的作用交换，我们可以考虑范畴
\[
  \Whit(G)\rd_{glob} :=
  D(\overline{\Bun}\rd_N)^{(\mathcal H_N,\chi),\mathcal H_{G,N}}
  = D(\Bun^{\mathrm{rat},\mathrm{red}}_N)^{(\mathcal H_N,\chi)},
\]
它是 $D(\overline{\Bun}\rd_N)^{\mathcal H_{G,N}}=D(\Bun^{\mathrm{rat},\mathrm{red}}_N)$ 的满子范畴。

我们有伴随函子对
\[
  \ind_{\Whit\rd,glob} : D(\Bun^{\mathrm{rat},\mathrm{red}}_N)
  \rightleftarrows \Whit(G)\rd_{glob} : \oblv_{\Whit\rd,glob}.
\]

\ob{6.2.1}{blk:6.2.1}\emph{典范对象。}范畴 $\Whit(G)\rd_{glob}$ 有典范的单位对象，记作 $\one_{\Whit(G)\rd_{glob}}$：它由把 $\ind\rd_N$ 作用于 $D(\overline{\Bun}\rd_N)$ 中某个典范的 $(\mathcal H_N,\chi)$-等变对象而得。

所说的对象是如下特征层的 $!$-（同时也是 $\ast$-）延拓：该特征层定义在 $\Bun\rd_N\subset\overline{\Bun}\rd_N$ 上，其中 $\Bun\rd_N$ 是 $\overline{\Bun}\rd_N$ 中 Plücker 数据 $\kappa$ 既无零点也无极点、且 $\gamma_i\rd$ 均为同构的那个轨迹。特征层由沿留数映射 $\Bun\rd_N\to\Ga$ 拉回 Artin-Schreier 层而得。

\ob{6.3}{blk:6.3}\textbf{$\Whit(G)_{glob}$ 与 $\Whit(G)\rd_{glob}$ 的关系。}

在 [GL:Ran] 中我们将证明：

\begin{othm}{命题}{6.3.1}{prop:6.3.1}
范畴 $\Whit(G)_{glob}$ 带有（刚性）幺半范畴 $D(\Gr_{Z_G,\mathrm{Ran}_X,un})$ 的自然作用，且有典范等价
\[
  \Whit(G)\rd_{glob}\simeq
  \Whit(G)_{glob}\underset{D(\Gr_{Z_G,\mathrm{Ran}_X,un})}{\otimes} D(\Bun_{Z_G}).
\]
\end{othm}

特别地，有伴随函子对
\[
  \Whit(G)_{glob}\rightleftarrows\Whit(G)\rd_{glob},
\]
其右伴随全忠实。

注意，$\one_{\Whit(G)_{glob}}$ 在上述函子 $\Whit(G)_{glob}\to\Whit(G)\rd_{glob}$ 下的像，正是上面那个对象 $\one_{\Whit(G)\rd_{glob}}$。

\ob{6.3.2}{blk:6.3.2}上述命题与定理 5.3.3 结合，给出 $\Whit(G)\rd_{glob}$ 的如下谱描述：

\begin{othm}{推论}{6.3.3}{cor:6.3.3}
我们有
\[
  \Whit(G)\rd_{glob}\simeq
  \Rep(\Ch G,\mathrm{Ran}_X,un)
  \underset{\Rep(\Ch G/[\Ch G,\Ch G],\mathrm{Ran}_X,un)}{\otimes}
  \QCoh\bigl(\LocSys_{\Ch G/[\Ch G,\Ch G]}\bigr).
\]
\end{othm}

\begin{annote}
命题 6.3.1 与推论 6.3.3 是同一件事的两侧。几何侧说：既约 Whittaker 范畴是把 $\Whit(G)_{glob}$ 沿 $\Gr_{Z_G}$ 的（Hecke）作用「张量缩小」到 $D(\Bun_{Z_G})$；谱侧说：这相当于把 $\Rep(\Ch G)$ 换成它商去换位子群 $\Ch G/[\Ch G,\Ch G]$ 后的表示，并把 $\QCoh(\LocSys_{\Ch G})$ 换成 $\Ch G/[\Ch G,\Ch G]$ 的局部系统叠上的拟凝聚层。两边都体现「$Z_G$ 连通」这一假设的作用：它保证中心部分可以在两侧一致地剥离。\par\vspace{0.15em}{\raggedleft 【来源：原书 p.11 §6.3--6.3.2（A）】\par}
\end{annote}

\section{扩展 Whittaker 范畴}
\markboth{扩展 Whittaker 范畴}{}

现在转向本文的主角，扩展 Whittaker 范畴，记作 $\Whit(G)\ext_{glob}$。

Whittaker 范畴 $\Whit(G)\rd_{glob}$ 理应从 $D(\Bun_G)$ 接受一个函子，它「抓住」自守 D-模的全部非退化 Fourier 系数（在 $G$ 的中心连通的假设下）。扩展 Whittaker 范畴 $\Whit(G)\ext_{glob}$ 则应当抓住全部 Fourier 系数。

让我们先解释函数论上的类比。令 $\Char^{\mathrm{n.d.}}$ 记 $N(\A)$ 上在 $N(K)$ 上平凡的非退化特征所构成的空间。$T(\A)$ 在 $N(\A)$ 上的伴随作用给出 $T(K)/Z_G(K)$ 在 $\Char^{\mathrm{n.d.}}$ 上的作用，而这个作用是单可迁的（依据 $Z_G$ 连通的假设）。于是我们得到同构
\begin{equation}
  Z_G(K)N(K)\backslash G(\A)/G(\mathbb O)
  \simeq
  T(K)\backslash\Bigl(N(K)\backslash G(\A)/G(\mathbb O)\times\Char^{\mathrm{n.d.}}\Bigr).
  \label{eq:7.1}
\end{equation}

此外，在 $N(K)\backslash G(\A)/G(\mathbb O)\times\Char^{\mathrm{n.d.}}$ 上的函数空间上，可以施加自然的 $N(\A)$-等变条件（该条件依赖于 $\Char^{\mathrm{n.d.}}$ 的一个点），而这个条件在 $T(K)$-作用下不变。所得的函数空间因而可以等同于 $Z_G(K)N(K)\backslash G(\A)/G(\mathbb O)$ 上的 Whittaker 函数空间。

现在终于可以解释 $D(\Bun^{\mathrm{rat},\mathrm{red},\mathrm{ext}}_N)$ 及其满子范畴 $\Whit(G)\ext_{glob}$ 在函数论上的类比了。具体地说，令 $\Char$ 记 $N(\A)$ 上在 $N(K)$ 上平凡的全部特征所构成的空间。范畴 $D(\Bun^{\mathrm{rat},\mathrm{red},\mathrm{ext}}_N)$ 是
\[
  T(K)\backslash\Bigl(N(K)\backslash G(\A)/G(\mathbb O)\times\Char\Bigr)
\]
上函数空间的几何类比；而 $\Whit(G)\ext_{glob}$ 是其中「要求关于 $N(\A)$ 按依赖于 $\Char$ 的点的特征作等变」所得函数子空间的几何类比。

特别地要注意：当我们特殊到 $\{0\}\subset\Char$ 时，就回到
\[
  T(K)N(\A)\backslash G(\A)/G(\mathbb O)
\]
上的函数集合。

\begin{annote}
$\Char$ 与 $\Char^{\mathrm{n.d.}}$ 的区别是本节的关键：$\Char^{\mathrm{n.d.}}$ 只是非退化特征那一片，$T(K)$ 在其上单可迁地作用，故 (7.1) 把「$Z_G$ 处的 Whittaker 函数空间」等同于「$T(K)$ 作用下的非退化特征丛的截面」；而 $\Char$ 包含退化特征，其上的等变条件会退化，因而得到的是「全部 Fourier 系数」而非只有非退化部分。$D(\Bun^{\mathrm{rat},\mathrm{red},\mathrm{ext}}_N)$ 正是 $\Char$ 全体的几何对应物。\par\vspace{0.15em}{\raggedleft 【来源：原书 p.12 §7（A）；经典 Whittaker 模型中「非退化特征」与「全部特征」的对比（A）】\par}
\end{annote}

\ob{7.1}{blk:7.1}\textbf{范畴。}空间 $\overline{\Bun}^{\mathrm{red},\mathrm{ext}}_N$ 与范畴
\[
  D(\Bun^{\mathrm{rat},\mathrm{red},\mathrm{ext}}_N)
  := D\bigl(\overline{\Bun}^{\mathrm{red},\mathrm{ext}}_N\bigr)^{\mathcal H_{G,N}}
\]
的定义，与 $\overline{\Bun}\rd_N$ 和 $D(\Bun^{\mathrm{rat},\mathrm{red}}_N)^{\mathcal H_{G,N}}$ 的定义逐字相同，唯一的区别是：现在映射
\[
  \gamma_i^{\mathrm{red},\mathrm{ext}} : \Ch\alpha_i(P_T)\to\omega_X
\]
允许取零。

与上面一样，我们有伴随函子对
\[
  \ind_N^{\mathrm{red},\mathrm{ext}} :
  D\bigl(\overline{\Bun}^{\mathrm{red},\mathrm{ext}}_N\bigr)
  \rightleftarrows
  D(\Bun^{\mathrm{rat},\mathrm{red},\mathrm{ext}}_N) : \oblv_N^{\mathrm{red},\mathrm{ext}}.
\]

先前的讨论对群胚 $\mathcal H_N$ 同样适用，于是得到相应的满子范畴
\[
  \Whit(G)\ext_{glob} :=
  D\bigl(\overline{\Bun}^{\mathrm{red},\mathrm{ext}}_N\bigr)^{(\mathcal H_N,\chi),\mathcal H_{G,N}}
  \subset D\bigl(\overline{\Bun}^{\mathrm{red},\mathrm{ext}}_N\bigr)^{\mathcal H_{G,N}}
  =: D(\Bun^{\mathrm{rat},\mathrm{red},\mathrm{ext}}_N).
\]

我们有相应的伴随函子对
\[
  \ind_{\Whit\ext_{glob}} : D(\Bun^{\mathrm{rat},\mathrm{red},\mathrm{ext}}_N)
  \rightleftarrows \Whit(G)\ext_{glob} : \oblv_{\Whit\ext},
\]
其右伴随全忠实。

\ob{7.2}{blk:7.2}\textbf{扩展 Whittaker 范畴与 $D(\Bun\rt_T)$。}注意，把 $\gamma_i\rd$ 的资料置零，$\bBun_B$ 就闭嵌入到 $\overline{\Bun}^{\mathrm{red},\mathrm{ext}}_N$ 中。把这个映射记作 $\fI$。

我们把 $\fI$ 看作诱导出「闭嵌入」$\fI\rt : \Bun\rt_B\hookrightarrow \Bun^{\mathrm{rat},\mathrm{red},\mathrm{ext}}_N$，意思是说，我们有伴随函子对
\[
  (\fI\rt)_! := (\fI\rt)_* :
  D(\Bun\rt_B)\rightleftarrows D(\Bun^{\mathrm{rat},\mathrm{red},\mathrm{ext}}_N) : (\fI\rt)^!,
\]
其中 $(\fI\rt)_!$ 全忠实，并且使下图交换：
\[
\begin{tikzcd}[row sep=2.9em, column sep=3.4em, ampersand replacement=\&]
  D(\Bun\rt_B) \arrow[r] \arrow[d, "\oblv_B"'] \&
  D(\Bun^{\mathrm{rat},\mathrm{red},\mathrm{ext}}_N) \arrow[d, "\oblv_N^{\mathrm{red},\mathrm{ext}}"] \\
  D(\bBun_B) \arrow[r] \&
  D(\overline{\Bun}^{\mathrm{red},\mathrm{ext}}_N)
\end{tikzcd}
\]
（以及把相应函子换成其伴随后所得的另外三个图）。

上述伴随对把相应的 $\mathcal H_N$-等变范畴互相对应，从而诱导出伴随对
\[
  (\fI\rt)_! := (\fI\rt)_* :
  D(\Bun\rt_T)\rightleftarrows \Whit(G)\ext_{glob} : (\fI\rt)^!.
\]

\ob{7.3}{blk:7.3}\textbf{扩展与非扩展的范畴。}我们有开嵌入
\[
  \fJ : \overline{\Bun}\rd_N\hookrightarrow\overline{\Bun}^{\mathrm{red},\mathrm{ext}}_N .
\]
我们把 $\fJ$ 看作诱导出「开嵌入」
\[
  \fJ\rt : \Bun^{\mathrm{rat},\mathrm{red}}_N\hookrightarrow\Bun^{\mathrm{rat},\mathrm{red},\mathrm{ext}}_N,
\]
它对应的是所有 $\gamma_i\rd$ 均非零的那个轨迹。

严格地说，这意味着我们有如下限制函子
\[
  (\fJ\rt)^! := (\fJ\rt)^* :
  D(\Bun^{\mathrm{rat},\mathrm{red},\mathrm{ext}}_N)\to D(\Bun^{\mathrm{rat},\mathrm{red}}_N),
\]
它有右伴随，记作 $(\fJ\rt)_*$，且该右伴随全忠实。

这些函子就函子
\[
  \ind_{N^{\mathrm{red},\mathrm{ext}}} : D(\overline{\Bun}^{\mathrm{red},\mathrm{ext}}_N)
  \rightleftarrows D(\Bun^{\mathrm{rat},\mathrm{red},\mathrm{ext}}_N) : \oblv_{N^{\mathrm{red},\mathrm{ext}}}
\]
与
\[
  \ind_{N\rd} : D(\overline{\Bun}\rd_N)
  \rightleftarrows D(\Bun^{\mathrm{rat},\mathrm{red}}_N) : \oblv_{N\rd}
\]
而言，给出四个交换图。

\ob{7.3.1}{blk:7.3.1}上述函子对把相应的 Whittaker 子范畴互相对应，于是我们得到伴随对
\[
  (\fJ\rt)^* : \Whit(G)\rd_{glob}\rightleftarrows\Whit(G)\ext_{glob} : (\fJ\rt)_*,
\]
其中 $(\fJ\rt)_*$ 全忠实。

下面我们来证明：

\begin{othm}{命题}{7.3.2}{prop:7.3.2}
$(\fJ\rt)^!$ 的部分定义的左伴随（记作 $(\fJ\rt)_!$）在子范畴 $\Whit(G)\rd_{glob}\subset D(\Bun^{\mathrm{rat},\mathrm{red}}_N)$ 上有定义。
\end{othm}

由命题形式地推出：$(\fJ\rt)_!$ 把 $\Whit(G)\ext_{glob}$ 送到 $\Whit(G)\rd_{glob}$，并且全忠实，从而是 $(\fJ\rt)^!$ 的右逆。

\begin{proof}
把 Hecke 范畴看作 $\Rep(\Ch G,\mathrm{Ran}_X,un)$ 的 Ran 版本，它作用在所有涉及的范畴上。这个作用与函子 $(\fJ\rt)^!$ 交换；又因 $\Rep(\Ch G,\mathrm{Ran}_X,un)$ 是刚性的，它也与 $(\fJ\rt)_!$ 在后者有定义的子范畴上交换。

由定理 5.3.3 推出，$\Whit(G)\rd_{glob}$ 在 Hecke 作用下由单个对象生成，即由 $\one_{\Whit(G)\rd_{glob}}$ 生成。因此只需证明
\[
  (\fJ\rt)_!\bigl(\one_{\Whit(G)\rd_{glob}}\bigr)
\]
有定义。而后者显然如此：$\one_{\Whit(G)\rd_{glob}}$ 是把 $\ind_N^{\mathrm{red},\mathrm{ext}}$ 作用于 $D(\overline{\Bun}^{\mathrm{red},\mathrm{ext}}_N)$ 中一个完整（holonomic）对象而得。
\end{proof}

\begin{othm}{拟命题}{7.3.3}{prop:7.3.3}
函子 $(\fI\rt)_* : D(\Bun\rt_T)\to\Whit(G)\ext_{glob}$ 有左伴随。
\end{othm}

我们只在 $G$ 的半单秩为 $1$ 时证明本命题：

\begin{proof}
对对象 $\mathcal F\in\Whit(G)\ext_{glob}$，考虑好三角
\[
  (\fJ\rt)_!\circ(\fJ\rt)^!(\mathcal F)\to\mathcal F\to\mathcal F_1,
\]
其中 $\mathcal F_1$ 满足 $(\fJ\rt)^!(\mathcal F_1)=0$。

然而，当 $G$ 的半单秩为 $1$ 时，后一条件意味着 $\mathcal F_1$ 落在 $(\fI\rt)_*$ 的本质像中，于是函子 $(\fI\rt)^*$ 在它上面有定义。
\end{proof}

\begin{annote}
两处省略需要点破。7.3.2 的证明里，「在 Hecke 作用下由单个对象生成」用的是定理 5.3.3 把 $\Whit(G)\rd_{glob}$ 与 $\Rep(\Ch G,\mathrm{Ran}_X,un)$ 等同起来，而后者作为刚性幺半范畴在其正则对象下生成自身；「刚性」是为了让左伴随与 Hecke 作用交换。7.3.3 的证明里，半单秩 $1$ 只有两个分层（$\fI$ 的像与 $\fJ$ 的像），故 $\mathcal F_1$ 必落在 $\fI$ 的像上，$(\fI\rt)^*$ 才有定义。\par\vspace{0.15em}{\raggedleft 【来源：原书 p.14 §7.3.2--7.3.3（A）；译者按（用两个分层的情形补足原文省略的推理）】\par}
\end{annote}

\ob{7.4}{blk:7.4}\textbf{系数函子。}我们有自然的遗忘映射
\[
  \fR : \overline{\Bun}^{\mathrm{red},\mathrm{ext}}_N\to\bBun_B .
\]
拉回诱导函子
\[
  \fR^! : D(\bBun_B)\to D(\overline{\Bun}^{\mathrm{red},\mathrm{ext}}_N).
\]
我们把 $\fR$ 看作诱导出「叠」的映射 $\Bun^{\mathrm{rat},\mathrm{red},\mathrm{ext}}_N\to\Bun\rt_B$。严格地说，我们的意思是存在典范定义的函子
\[
  (\fR\rt)^! : D(\Bun\rt_B)\to D(\Bun^{\mathrm{rat},\mathrm{red},\mathrm{ext}}_N),
\]
它使下图交换：
\[
\begin{tikzcd}[row sep=2.9em, column sep=4.2em, ampersand replacement=\&]
  D(\Bun\rt_B) \arrow[r, "(\fR\rt)^!"] \arrow[d, "\oblv_B"'] \&
  D(\Bun^{\mathrm{rat},\mathrm{red},\mathrm{ext}}_N) \arrow[d, "\oblv_N^{\mathrm{red},\mathrm{ext}}"] \\
  D(\bBun_B) \arrow[r, "\fR^!"'] \&
  D(\overline{\Bun}^{\mathrm{red},\mathrm{ext}}_N)
\end{tikzcd}
\]
并且把纵向箭头分别换成 $\ind_B$ 与 $\ind_N^{\mathrm{red},\mathrm{ext}}$ 后的另一个图也交换。

我们定义函子
\[
  \coeff\ext : D(\Bun_G)\to\Whit(G)\ext_{glob}
\]
为复合
\[
  \coeff\ext := \ind_{\Whit\ext,glob}\circ(\fR\rt)^!\circ(\fP\rt)^! .
\]

复合
\[
  \coeff\rd := (\fJ\rt)^!\circ\coeff\ext : D(\Bun_G)\to\Whit(G)\rd_{glob}
\]
是通常的 Whittaker 系数函子。

注意，复合
\[
  (\fI\rt)^!\circ\coeff\ext : D(\Bun_G)\to D(\Bun\rt_T)
\]
与函子 $\CT\rt$ 一致。

\begin{othm}{拟命题}{7.4.1}{prop:7.4.1}
函子 $\coeff\ext$ 有左伴随。
\end{othm}

我们把这一左伴随函子记作 $\Poinc\ext_!$。复合
\[
  \Poinc\rd_! := \Poinc\ext_!\circ(\fJ\rt)_! :
  \Whit(G)\rd_{glob}\to D(\Bun_G)
\]
是 $\coeff\rd$ 的左伴随，即通常的 Poincaré 级数函子。

下面我们只在 $G$ 的半单秩为 $1$ 时证明本定理。

\begin{proof}
如同定理 7.3.3 的证明，只需分别证明 $\Poinc\rd_!$ 有定义，以及 $\CT\rt$ 的左伴随有定义。后者已在 §3.2 中讨论过。前者如命题 7.3.2 的证明。
\end{proof}

\begin{annote}
原文此处称「定理 7.3.3」，但 7.3.3 是拟命题，7.3.2 才是命题；译文按原样译出并在工程文件疑点清单中记录。\par\vspace{0.15em}{\raggedleft 【来源：原书 p.15 §7.4（A）】\par}
\end{annote}

\ob{7.5}{blk:7.5}\textbf{全忠实性。}我们处理 $GL_n$ 情形下几何朗兰兹的方法，关键在于下述论断：

\begin{othm}{拟定理}{7.5.1}{thm:7.5.1}
当 $G=GL_n$ 时，函子 $\coeff\ext$ 全忠实。
\end{othm}

对其他的约化群，我们不知道类似的论断是否成立。

\ob{7.5.2}{blk:7.5.2}下面朝定理 7.5.1 的证明走一步。

对约化群 $G$ 与抛物子群 $P$，可以仿照 $D(\Bun\rt_B)$（用部分的 Plücker 数据）定义范畴 $D(\Bun\rt_P)$。我们有伴随函子对
\[
  (\fP\rt_P)_! : D(\Bun\rt_P)\rightleftarrows D(\Bun_G) : (\fP\rt_P)^!,
\]
并且如同定理 1.3.1，我们有：

\begin{othm}{拟定理}{7.5.3}{thm:7.5.3}
函子 $(\fP\rt_P)^!$ 全忠实。
\end{othm}

现在取 $G=GL_n$，$P$ 为抛物子群 $(n-1,1)$。下述论断「应当」可以用经典的迭代 Fourier 变换手续证明：

\begin{othm}{拟定理}{7.5.4}{thm:7.5.4}
存在范畴的等价
\[
  D(\Bun\rt_P)\simeq\Whit(G)\ext_{glob}.
\]
\end{othm}

复合
\[
  D(\Bun_G)\xrightarrow{(\fP\rt_P)^!} D(\Bun\rt_P)\to\Whit(G)\ext_{glob}
\]
就是函子 $\coeff\ext$。

显然，上述两个定理蕴含定理 7.5.1。

\begin{annote}
这三步的用意是：把「$G$ 的扩展 Whittaker 范畴」换成「抛物 $P$ 的有理范畴」，而后者是更接近 $D(\Bun_G)$ 的对象。拟定理 7.5.3 说 $(\fP\rt_P)^!$ 全忠实，拟定理 7.5.4 说 $D(\Bun\rt_P)$ 与扩展 Whittaker 范畴等价；两者一夹，就得到拟定理 7.5.1。原文把 7.5.1、7.5.3、7.5.4 都标为拟定理，后文却简称「定理」，译文照录其简称。\par\vspace{0.15em}{\raggedleft 【来源：原书 p.16 §7.5.2（A）；译者按（把三步的逻辑关系理清）】\par}
\end{annote}

\section{粘合扩展 Whittaker 范畴，以及谱图景}
\markboth{粘合扩展 Whittaker 范畴，以及谱图景}{}

本节的材料没有经典类比。

我们想描述范畴 $\Whit(G)\ext_{glob}$ 如何由它的子范畴粘合而成，例如 $\Whit(G)\rd_{glob}$ 与 $D(\Bun\rt_T)$，这些子范畴对应 $\Bun^{\mathrm{rat},\mathrm{red},\mathrm{ext}}_N$ 的各个「分层」，而分层由哪些 $\gamma_i\rd$ 取零的模式决定。

粘合数据由这些子范畴之间的函子给出，例如
\[
  (\fI\rt)^!\circ(\fJ\rt)_! : \Whit(G)\rd_{glob}\to D(\overline{\Bun}\rt_T).
\]

\begin{annote}
原文此处把靶写成 $D(\overline{\Bun}\rt_T)$（$\Bun$ 上有横线、同时有 $rat$ 上标）。$(\fI\rt)^!$ 的靶应是 $D(\Bun\rt_T)$（横线表示紧化，与 $rat$ 不能并存）；由 §7.2 的伴随对与 §8.1 的 8.1 式，靶应为 $D(\Bun\rt_T)$。译文照录原文，见工程文件疑点清单。\par\vspace{0.15em}{\raggedleft 【来源：原书 p.16 §8（A）；译者按（类型检查）】\par}
\end{annote}

\ob{8.1}{blk:8.1}\textbf{半单秩 $1$ 的粘合数据。}让我们注意：当 $G$ 的半单秩为 $1$ 时，$\Bun^{\mathrm{rat},\mathrm{red},\mathrm{ext}}_N$ 中「只有两个分层」，即 $\Bun\rt_B$ 与 $\Bun^{\mathrm{rat},\mathrm{red},\mathrm{ext}}_N$，因此范畴 $\Whit(G)\ext_{glob}$ 完全由函子 $(\fI\rt)^!\circ(\fJ\rt)_!$ 描述。

事实上，$D(\Bun^{\mathrm{rat},\mathrm{red},\mathrm{ext}}_N)$ 可以等同于这样的范畴，其对象是三元组 $(\mathcal F_1,\mathcal F_2,\alpha)$，其中
\[
  \mathcal F_1\in\Whit(G)\rd_{glob},\qquad
  \mathcal F_2\in D(\Bun\rt_T),
\]
而
\[
  \alpha\in\Hom_{D(\Bun\rt_T)}\bigl((\fI\rt)^!\circ(\fJ\rt)_!(\mathcal F_1),\mathcal F_2\bigr).
\]
态射按自然的方式取为三元组之间的态射。

对半单秩更高的群，范畴 $\Whit(G)\ext_{glob}$ 由一族这样的函子之间的图来描述，这些函子对应所有互相包含的抛物子群对。

\begin{annote}
「粘合」在这里是「分层—粘合」意义上的：把 $\Bun^{\mathrm{rat},\mathrm{red},\mathrm{ext}}_N$ 按 $\gamma_i\rd$ 取零的模式分成若干层，每层的范畴已知（$\Whit(G)\rd_{glob}$ 对应全非零层，$D(\Bun\rt_T)$ 对应全零层），层与层之间的函子 $(\fI\rt)^!\circ(\fJ\rt)_!$ 就是粘合数据。半单秩 $1$ 时只有两层，故三元组的描述完整；秩高时层数增多，要用「抛物对之间的函子图」。\par\vspace{0.15em}{\raggedleft 【来源：原书 p.16 §8.1（A）；译者按（把分层的含义说清）】\par}
\end{annote}

\ob{8.2}{blk:8.2}\textbf{谱描述。}由推论 6.3.3，我们有
\[
  \Whit(G)\rd_{glob}\simeq
  \Rep(\Ch G,\mathrm{Ran}_X,un)
  \underset{\Rep(\Ch G/[\Ch G,\Ch G],\mathrm{Ran}_X,un)}{\otimes}
  \QCoh\bigl(\LocSys_{\Ch G/[\Ch G,\Ch G]}\bigr),
\]
而由定理 4.6.1，我们有
\[
  D(\Bun\rt_T)\simeq T_{\LocSys_{\Ch B}/\LocSys_{\Ch G}}\text{-}\mathrm{mod}^r .
\]
我们当前的目标是在这些项下把函子 $(\fI\rt)^!\circ(\fJ\rt)_!$ 认出来，即把它认成函子
\begin{equation}
\begin{aligned}
&\Rep(\Ch G,\mathrm{Ran}_X,un)
  \underset{\Rep(\Ch G/[\Ch G,\Ch G],\mathrm{Ran}_X,un)}{\otimes}
  \QCoh\bigl(\LocSys_{\Ch G/[\Ch G,\Ch G]}\bigr)\\
&\qquad\longrightarrow
  T_{\LocSys_{\Ch B}/\LocSys_{\Ch G}}\text{-}\mathrm{mod}^r .
\end{aligned}
\label{eq:8.1}
\end{equation}

下述论断受 V. Drinfeld 提出、A. Arinkin 证明的「Gap 定理」启发。

\begin{othm}{拟定理}{8.2.1}{thm:8.2.1}
粘合函子 \eqref{eq:8.1} 典范同构于复合
\[
\begin{aligned}
&\Rep(\Ch G,\mathrm{Ran}_X,un)
  \underset{\Rep(\Ch G/[\Ch G,\Ch G],\mathrm{Ran}_X,un)}{\otimes}
  \QCoh\bigl(\LocSys_{\Ch G/[\Ch G,\Ch G]}\bigr)\\
&\qquad\xrightarrow{\Loc_{\Ch G,spec}}
  \QCoh(\LocSys_{\Ch G})
  \xrightarrow{(\fP\st)^!}
  T_{\LocSys_{\Ch B}/\LocSys_{\Ch G}}\text{-}\mathrm{mod}^r .
\end{aligned}
\]
\end{othm}

在 [GL:funct] 中我们将说明，定理 8.2.1 如何允许我们对 $G=GL_2$ 构造一个单方向的函子
\[
  \Psi_G : \QCoh_N(\LocSys_{\Ch G})\to D(\Bun_G).
\]

\begin{annote}
Gap 定理（间隙定理）是 Drinfeld 提出、Arinkin 证明的关于 $\Bun_G$ 上 D-模与谱侧之间「间隙」的定理；原文只说本文的论断受其启发，并未用其内容。这里的关键是：粘合函子 $(\fI\rt)^!\circ(\fJ\rt)_!$ 本来是用几何方式定义的，而 \eqref{eq:8.1} 把它翻译到谱侧，认成「先 $\Loc_{\Ch G,spec}$、再 $(\fP\st)^!$」。这一步是 §8 的全部内容，也是 [GL:funct] 构造 $\Psi_G$ 的出发点。\par\vspace{0.15em}{\raggedleft 【来源：原书 p.17 §8.2（A）；译者按（把该论断在全文中的位置说清）】\par}
\end{annote}


\backmatter
\chapter*{文献}
\markboth{文献}{}
\addcontentsline{toc}{chapter}{文献}

\newcommand{\refitem}[2]{%
  \par\noindent\hangindent=5.6em\hangafter=1
  \makebox[5.6em][l]{\texttt{#1}}#2\par\smallskip}

\refitem{[BD]}{Beilinson, A., Drinfeld, V., \emph{Chiral algebras}, American Mathematical Society Colloquium Publications, 51, 2004.}
\refitem{[DrGa1]}{Drinfeld, V., Gaitsgory, D., \emph{The category of D-modules on $\Bun_G$}, 预印本.}
\refitem{[Ga1]}{Gaitsgory, D., \emph{The vanishing conjecture}, 预印本.}
\refitem{[Ga2]}{Gaitsgory, D., \emph{Twisted Whittaker model and factorizable sheaves}, 预印本.}
\refitem{[GL:DG]}{Gaitsgory, D., \emph{Notes on Geometric Langlands: Generalities on DG categories}, 预印本.}
\refitem{[GL:formulation]}{Gaitsgory, D., \emph{Notes on Geometric Langlands: Formulation of the GL conjecture}, 预印本.}
\refitem{[GL:contr]}{Gaitsgory, D., \emph{Notes on Geometric Langlands: Contractibility of the space of rational maps}, 预印本.}
\refitem{[GL:Ran]}{Gaitsgory, D., \emph{Notes on Geometric Langlands: Categories over the Ran space}, 预印本.}
\refitem{[GL:Eis]}{Gaitsgory, D., \emph{Notes on Geometric Langlands: What acts on Eisenstein series}, 预印本.}
\refitem{[GL:funct]}{Gaitsgory, D., \emph{Notes on Geometric Langlands: Construction of the functor}, 预印本.}

\medskip
\noindent\emph{说明。}原书文献只有代号、作者与篇名，未给出期刊、卷页或 arXiv 编号，译文照录，不代补。系列名 \emph{Notes on Geometric Langlands} 的十册之间以代号互引，正文中的 \texttt{[GL:…]} 即指此表。


\end{document}
